% ---------------------------------------------------------------------------
% Building High-Performance MCP Applications - preamble
% Every visual decision in the book lives here. Chapters carry no formatting.
% ---------------------------------------------------------------------------

\usepackage[T1]{fontenc}
\usepackage[utf8]{inputenc}
\usepackage{lmodern}

% Prefer Libertinus for text/math and Inconsolata for code; fall back quietly
% so the book still builds on a minimal TeX Live.
\IfFileExists{libertinus.sty}{%
  \usepackage[mono=false]{libertinus}%
}{%
  \IfFileExists{newtxtext.sty}{\usepackage{newtxtext}\usepackage{newtxmath}}{}%
}
\IfFileExists{inconsolata.sty}{\usepackage[varqu,varl,scaled=0.95]{inconsolata}}{}

\usepackage{microtype}
\usepackage[english]{babel}
\usepackage{csquotes}

% ---------------------------------------------------------------------------
% Page geometry. Generous outer margin: this book has a lot of marginal
% asides and wide code, and a cramped measure makes both miserable.
% ---------------------------------------------------------------------------
\usepackage[
  paperwidth=7.5in, paperheight=9.25in,
  inner=0.95in, outer=0.85in,
  top=0.9in, bottom=1.05in,
  marginparwidth=0in, marginparsep=0pt,
  headsep=18pt, footskip=32pt,
]{geometry}

\usepackage{graphicx}
\usepackage{xcolor}
\usepackage{booktabs}
\usepackage{longtable}
\usepackage{array}
\usepackage{tabularx}
\usepackage{multirow}
\usepackage{enumitem}
\usepackage{amsmath}
\usepackage{amssymb}
\usepackage{fancyhdr}
\usepackage{titlesec}
\usepackage{titletoc}
\usepackage{needspace}
\usepackage{ragged2e}
\usepackage[most]{tcolorbox}
\usepackage{listings}
\usepackage{caption}
\usepackage{float}
\usepackage[hidelinks,bookmarksnumbered]{hyperref}

\setlist{itemsep=2pt, parsep=0pt, topsep=5pt}
\setlength{\parskip}{0pt}
\setlength{\emergencystretch}{3em}
\widowpenalty=10000
\clubpenalty=10000

% ---------------------------------------------------------------------------
% Palette. Ink is not pure black; pure black on cream paper reads harsh.
% ---------------------------------------------------------------------------
\definecolor{ink}{HTML}{15181D}
\definecolor{muted}{HTML}{5B6472}
\definecolor{rule}{HTML}{C9CED6}
\definecolor{wash}{HTML}{F4F5F7}
\definecolor{accent}{HTML}{0F5C8C}      % protocol blue
\definecolor{accentsoft}{HTML}{E7F0F6}
\definecolor{warm}{HTML}{B4531A}        % legacy / deprecation orange
\definecolor{warmsoft}{HTML}{FBF0E7}
\definecolor{good}{HTML}{2C6E49}        % measurements green
\definecolor{goodsoft}{HTML}{E9F2EC}
\definecolor{danger}{HTML}{9B2226}      % security red
\definecolor{dangersoft}{HTML}{FAECEC}

\color{ink}

% ---------------------------------------------------------------------------
% Headings
% ---------------------------------------------------------------------------
\titleformat{\part}[display]
  {\normalfont\sffamily\filcenter}
  {\color{muted}\large\scshape Part \thepart}
  {12pt}
  {\Huge\bfseries\color{ink}}
  [\vspace{6pt}{\color{rule}\rule{0.35\textwidth}{0.6pt}}]

\titleformat{\chapter}[display]
  {\normalfont\sffamily\raggedright}
  {\color{accent}\Large\bfseries\chaptertitlename\ \thechapter}
  {6pt}
  {\Huge\bfseries\color{ink}}
  [\vspace{4pt}{\color{accent}\rule{\textwidth}{1.2pt}}]
\titlespacing*{\chapter}{0pt}{6pt}{26pt}

\titleformat{\section}
  {\normalfont\sffamily\Large\bfseries\color{ink}}{\thesection}{0.7em}{}
\titleformat{\subsection}
  {\normalfont\sffamily\large\bfseries\color{ink}}{\thesubsection}{0.6em}{}
\titleformat{\subsubsection}
  {\normalfont\sffamily\normalsize\bfseries\color{muted}}{}{0em}{}
\titlespacing*{\section}{0pt}{18pt}{6pt}
\titlespacing*{\subsection}{0pt}{14pt}{4pt}
\titlespacing*{\subsubsection}{0pt}{12pt}{3pt}

\setcounter{secnumdepth}{2}
\setcounter{tocdepth}{1}

% ---------------------------------------------------------------------------
% Running heads
% ---------------------------------------------------------------------------
\pagestyle{fancy}
\fancyhf{}
\renewcommand{\headrulewidth}{0.4pt}
\renewcommand{\footrulewidth}{0pt}
\fancyhead[LE]{\sffamily\footnotesize\color{muted}\leftmark}
\fancyhead[RO]{\sffamily\footnotesize\color{muted}\rightmark}
\fancyfoot[LE,RO]{\sffamily\small\color{muted}\thepage}
\renewcommand{\chaptermark}[1]{\markboth{\textsc{Chapter \thechapter}\quad #1}{}}
\renewcommand{\sectionmark}[1]{\markright{\textsc{\S\thesection}\quad #1}}
\fancypagestyle{plain}{\fancyhf{}\renewcommand{\headrulewidth}{0pt}%
  \fancyfoot[LE,RO]{\sffamily\small\color{muted}\thepage}}

% The blank verso pages `\cleardoublepage` inserts inherit the running head,
% which puts a rule and a folio on an otherwise empty page. Strip them.
\makeatletter
\def\cleardoublepage{\clearpage
  \if@twoside
    \ifodd\c@page\else
      \hbox{}\thispagestyle{empty}\newpage
      \if@twocolumn\hbox{}\newpage\fi
    \fi
  \fi}
\makeatother

% ---------------------------------------------------------------------------
% Code listings
% ---------------------------------------------------------------------------
\lstdefinelanguage{jsonrpc}{
  keywords={true,false,null},
  keywordstyle=\color{warm},
  sensitive=true,
  morestring=[b]",
  stringstyle=\color{accent},
  comment=[l]{//},
  commentstyle=\color{muted}\itshape,
}
\lstdefinelanguage{httpwire}{
  morekeywords={POST,GET,DELETE,PUT,HTTP,Content-Type,Accept,Authorization,
    MCP-Protocol-Version,Mcp-Method,Mcp-Name,Mcp-Param-Region,WWW-Authenticate,
    X-Accel-Buffering,Cache-Control},
  keywordstyle=\color{accent}\bfseries,
  sensitive=true,
  morestring=[b]",
  stringstyle=\color{ink},
  comment=[l]{\#},
  commentstyle=\color{muted}\itshape,
}
\lstdefinelanguage{typescriptx}{
  morekeywords={async,await,const,let,var,function,return,if,else,for,while,
    class,extends,implements,interface,type,import,export,from,new,this,
    try,catch,finally,throw,typeof,instanceof,null,undefined,true,false,
    string,number,boolean,void,Promise,readonly,private,public},
  keywordstyle=\color{warm}\bfseries,
  sensitive=true,
  morestring=[b]",
  morestring=[b]',
  morestring=[b]`,
  stringstyle=\color{good},
  morecomment=[l]{//},
  morecomment=[s]{/*}{*/},
  commentstyle=\color{muted}\itshape,
}

\lstdefinestyle{book}{
  basicstyle=\ttfamily\footnotesize\color{ink},
  numbers=none,
  frame=leftline,
  framesep=8pt,
  xleftmargin=12pt,
  rulecolor=\color{rule},
  framerule=1.6pt,
  backgroundcolor=\color{wash},
  showstringspaces=false,
  breaklines=true,
  breakatwhitespace=false,
  breakindent=12pt,
  postbreak=\mbox{\textcolor{muted}{$\hookrightarrow$}\space},
  columns=fullflexible,
  keepspaces=true,
  tabsize=2,
  upquote=true,
  aboveskip=12pt,
  belowskip=12pt,
  captionpos=b,
  escapeinside={(*@}{@*)},
}
\lstset{style=book}

\lstdefinestyle{wire}{style=book, language=jsonrpc,
  backgroundcolor=\color{accentsoft}, rulecolor=\color{accent}}
\lstdefinestyle{http}{style=book, language=httpwire,
  backgroundcolor=\color{accentsoft}, rulecolor=\color{accent}}
\lstdefinestyle{py}{style=book, language=Python}
\lstdefinestyle{ts}{style=book, language=typescriptx}

% Shorthand environments so chapters stay readable in source form.
\lstnewenvironment{wire}[1][]{\lstset{style=wire,#1}}{}
\lstnewenvironment{http}[1][]{\lstset{style=http,#1}}{}
\lstnewenvironment{py}[1][]{\lstset{style=py,#1}}{}
\lstnewenvironment{ts}[1][]{\lstset{style=ts,#1}}{}
\lstnewenvironment{sh}[1][]{\lstset{style=book,language=bash,#1}}{}
\lstnewenvironment{plain}[1][]{\lstset{style=book,language=,#1}}{}

\newcommand{\code}[1]{\texttt{\small #1}}

% ---------------------------------------------------------------------------
% Boxes
% ---------------------------------------------------------------------------
\newtcolorbox{legacybox}[1]{
  enhanced, breakable,
  colback=warmsoft, colframe=warm,
  boxrule=0pt, leftrule=2.5pt,
  arc=0pt, outer arc=0pt,
  left=10pt, right=10pt, top=9pt, bottom=9pt,
  fonttitle=\sffamily\bfseries\small,
  coltitle=warm,
  title={LEGACY\quad #1},
  attach boxed title to top left={xshift=10pt, yshift=-2pt},
  boxed title style={colback=warmsoft, boxrule=0pt, arc=0pt},
  top=16pt,
}

\newtcolorbox{measurebox}[1]{
  enhanced, breakable,
  colback=goodsoft, colframe=good,
  boxrule=0pt, leftrule=2.5pt,
  arc=0pt, outer arc=0pt,
  left=10pt, right=10pt, top=16pt, bottom=9pt,
  fonttitle=\sffamily\bfseries\small,
  coltitle=good,
  title={MEASURED\quad #1},
  attach boxed title to top left={xshift=10pt, yshift=-2pt},
  boxed title style={colback=goodsoft, boxrule=0pt, arc=0pt},
}

\newtcolorbox{dangerbox}[1]{
  enhanced, breakable,
  colback=dangersoft, colframe=danger,
  boxrule=0pt, leftrule=2.5pt,
  arc=0pt, outer arc=0pt,
  left=10pt, right=10pt, top=16pt, bottom=9pt,
  fonttitle=\sffamily\bfseries\small,
  coltitle=danger,
  title={THREAT\quad #1},
  attach boxed title to top left={xshift=10pt, yshift=-2pt},
  boxed title style={colback=dangersoft, boxrule=0pt, arc=0pt},
}

\newtcolorbox{notebox}[1]{
  enhanced, breakable,
  colback=wash, colframe=muted,
  boxrule=0pt, leftrule=2.5pt,
  arc=0pt, outer arc=0pt,
  left=10pt, right=10pt, top=16pt, bottom=9pt,
  fonttitle=\sffamily\bfseries\small,
  coltitle=muted,
  title={#1},
  attach boxed title to top left={xshift=10pt, yshift=-2pt},
  boxed title style={colback=wash, boxrule=0pt, arc=0pt},
}

% The one-line thesis of a section. Used sparingly, on purpose.
% The box must start in vertical mode. A stray \noindent here opens horizontal
% mode, and a breakable tcolorbox that then splits across a page raises
% "You can't use \prevdepth in horizontal mode" — a fatal error that names a
% page number and no file, and that only appears once the text above it grows
% enough to move the box onto a boundary.
\newcommand{\keyidea}[1]{%
  \par\vspace{8pt}%
  \begin{tcolorbox}[enhanced, breakable, colback=accentsoft,
    colframe=accent, boxrule=0pt, leftrule=3pt, arc=0pt, outer arc=0pt,
    left=11pt, right=11pt, top=9pt, bottom=9pt]
  \sffamily\itshape #1
  \end{tcolorbox}%
  \vspace{4pt}\par\noindent\ignorespaces
}

% ---------------------------------------------------------------------------
% Figures
% ---------------------------------------------------------------------------
\captionsetup{
  font={small,sf}, labelfont={bf,color=accent}, textfont={color=muted},
  labelsep=quad, justification=justified, singlelinecheck=false, skip=8pt,
}
\graphicspath{{figures/}{../figures/}}

% \bookfig{basename}{caption}{label}[width fraction]
\newcommand{\bookfig}[4][0.86]{%
  \begin{figure}[htbp]
    \centering
    \includegraphics[width=#1\textwidth]{#2}
    \caption{#3}\label{fig:#4}
  \end{figure}
}

\newcommand{\figref}[1]{Figure~\ref{fig:#1}}
\newcommand{\chapref}[1]{Chapter~\ref{ch:#1}}
\newcommand{\secref}[1]{\S\ref{sec:#1}}
\newcommand{\tabref}[1]{Table~\ref{tab:#1}}

% ---------------------------------------------------------------------------
% Tables
% ---------------------------------------------------------------------------
\renewcommand{\arraystretch}{1.22}
\newcolumntype{L}[1]{>{\RaggedRight\arraybackslash}p{#1}}
\newcolumntype{R}[1]{>{\RaggedLeft\arraybackslash}p{#1}}
\newcommand{\thead}[1]{\sffamily\bfseries\small #1}

% ---------------------------------------------------------------------------
% Table of contents
% ---------------------------------------------------------------------------
\titlecontents{part}[0pt]
  {\vspace{16pt}\sffamily\bfseries\large\color{accent}}
  {}{}{}
\titlecontents{chapter}[0pt]
  {\vspace{6pt}\sffamily}
  {\bfseries\contentslabel{1.6em}}{\hspace*{-1.6em}}
  {\bfseries\titlerule*[0.6pc]{.}\contentspage}
\titlecontents{section}[2.0em]
  {\sffamily\small\color{muted}}
  {\contentslabel{2.4em}}{}
  {\titlerule*[0.6pc]{.}\contentspage}

\hypersetup{
  linkcolor=accent, citecolor=accent, urlcolor=accent,
  pdftitle={Building High-Performance MCP Applications},
  pdfauthor={Krimler},
  pdfsubject={Model Context Protocol, revision 2026-07-28},
  pdfkeywords={MCP, Model Context Protocol, agents, LLM, performance},
}
